\subsection{序域的代数扩张}

设 K 为一个序域，且 f 为 K{[}X{]} 中的多项式。我们说 f 在 K 中变号，如果存在 K 中的两个元素 a 和 b，使得 \(f(a)\) \(f(b) < 0\) ；这时我们就说 f 在 a 与 b 之间变号。

命题 3.——设 K 是一个序域，f 是 K 上的不可约多项式，且 f 在 K 中的 a 与 b 之间变号。则扩张 \(E = K[X]/(f)\) 容许一个 K 的序扩张的结构。

我们通过对 f 的次数 n 进行归纳来论证。当 n = 1 时，证明是平凡的。假设结论对次数 \(\leqslant n - 1\) 的情形成立，我们用反证法证明次数为 n 的情形；由定理 1，我们因此假设一个形如

\(1 + \sum_{i}p_{i}f_{i}^{2}(\mathbf{X})\equiv 0 (\mathrm{mod} f(\mathbf{X}))\) ，其中 \(f_{i}\in \mathbf{K}[\mathbf{X}], p_{i}\in \mathbf{K}\) 且 \(p_i\geqslant 0\)

不失一般性，我们可以假设这些 \(f_{i}\) 的次数 \(\leqslant n - 1\) （IV，第 11 页，推论）。于是

\[
1 + \sum_ {i} p _ {i} f _ {i} ^ {2} (\mathbf {X}) = h (\mathbf {X}) f (\mathbf {X})
\]

，其中 \(h \neq 0\) 的次数至多为 n - 2。在上面的不等式中把 X 换成 a 和 b，我们看到 \(h(a) f(a) > 0\) 且 \(h(b) f(b) > 0\) 。由于根据假设 f 在 a 与 b 之间变号，我们得出 \(h(a) h(b) < 0\) 。于是对不可约因子 \(g(\mathbf{X})\) （\(h(\mathbf{X})\) 的因子之一）有类似的不等式：即 \(g(a)g(b) < 0\) 。但 \(1 + \sum_{i}p_{i}f_{i}^{2}(\mathbf{X})\equiv 0 (\mathrm{mod} g(\mathbf{X}))\) ，这表明域 \(\mathrm{K}[\mathrm{X}] / (g)\) 不能是 \(\mathbf{K}\) 的序扩张（定理 1），与归纳假设相矛盾。

注记。--- 存在序域 K 上的不可约多项式 f，它们在 K 中不变号，但 \(\mathrm{K}[\mathrm{X}]/(f)\) 却容许具有 K 的序扩张的结构（参见 VI，第 43 页，习题 26，c））。

为了应用前面的命题，我们需要以下结果：

命题 4. --- 设 K 是一个序域，并设 \(f \in \mathbf{K}[\mathbf{X}]\) 。在 K 中存在一个区间，在该区间的补集上， \(f\) 与其最高次项具有相同的符号。

我们可以立即将问题化归为多项式为首一多项式的情形；于是可以写出 \(f(x) = x^{n}(1 + a_{1}x^{-1} + \cdots + a_{n}x^{-n})\) ，对 \(x \neq 0\) 。设

\[
\mathbf {M} = \sup \left(1, | a _ {1} | + \dots + | a _ {n} |\right).
\]

对于 \(|x| > M\) 我们有 \(1 + a_{1}x^{-1} + \cdots + a_{n}x^{-n} > 0\) ，这就完成了命题的证明。

推论 1. --- 序域的奇数有限次扩张的每一个扩张都容许序扩张的结构。

这样的扩张，由于是单生成的（V，第 40 页，定理 1），同构于 K{[}X{]}/(f)，其中 f 是一个奇数次的不可约多项式。那么只需证明 f 在 K 中变号（命题 3），这由命题 4 立即得出。

推论 2. —— 若 a 是序域 K 的正元素，则多项式 \(X^{2}-a\) 的每个分裂域 E 都容许具有 K 的序扩张的结构。

若 a 是 K 中的平方，则结论是平凡的。否则，多项式 \(f(\mathbf{X}) = \mathbf{X}^{2} - a\) 是不可约的且变号，因为 \(f(0) < 0\) ，且对 K 的某个区间的补集中的 x，\(f(x)\) 与 \(x^{2}\) 同号，即为正。我们现在可以通过应用命题 3 来完成证明。

注：--- 当序域 K 包含 K 的正元素 \(a\) 的“平方根”（多项式 \(\mathbf{X}^2 - a\) 的根）时，记号 \(\sqrt{a}\) 通常保留给正的平方根。若 K 不包含 \(b\) 与 \(-b\) （\(a\) 在域 E 中的平方根），则后者可以以两种方式成为 K 的序扩张，每种方式都由另一个通过把 \(b\) 映为 \(-b\) 的 K-自同构诱导而来。选择这两个序中的哪一个，便确定了 \(\sqrt{a}\) ：它是元素 \(b\) 与 \(-b\) 中为正的那一个。

如果 \(\pmb{a}\) 与 \(a'\) 是 \(\mathbf{K}\) 的两个正元素，其平方根位于 \(\mathbf{K}\) ，则 \(\sqrt{aa'} = \sqrt{a}\sqrt{a'}\) ，这由 \(\sqrt{a}\) 的定义以及符号法则得出。

